Parametrický výpočet starobního důchodu v~závislosti na úrovni příjmů odhaluje zásadní asymetrii mezi skupinami pojištěnců: zaměstnanec s~průměrnou mzdou dosáhne výrazně vyššího důchodu než \ac{OSVČ} se srovnatelnými celkovými náklady zaměstnavatele, a~to zejména v~případě paušálního daňového režimu, kde je výměra důchodu stropována pevnou platbou.
Tato konstrukce má přímou vazbu na sociální dialog: nárůst podílu \ac{OSVČ} na zaměstnanosti (srov.~obr.~\ref{fig:self_employment_map}), částečně podnícený vyhýbáním se kolektivnímu vyjednávání ze strany zaměstnavatelů (tzv.\ švarcsystém), oslabuje příjmovou základnu důchodového systému.
Slabé kolektivní vyjednávání tak působí systémově na dvou úrovních: přímou cestou stlačuje mzdy zaměstnanců a~nepřímo, prostřednictvím rozšiřování fiktivní samostatné výdělečné činnosti, erozi pojistné základny, z~níž se důchody financují.
Reforma sociálního dialogu --- zejména posílení mechanismů proti švarcsystému a~rozšíření sektorových \ac{KS} --- by přispěla i~ke stabilitě průběžného důchodového systému.

Parametrický výpočet starobního důchodu v~závislosti na úrovni příjmů odhaluje zásadní asymetrii mezi skupinami pojištěnců: zaměstnanec s~průměrnou mzdou dosáhne výrazně vyššího důchodu než \ac{OSVČ} se srovnatelnými celkovými náklady zaměstnavatele, a~to zejména v~případě paušálního daňového režimu, kde je výměra důchodu stropována pevnou platbou.
Tato konstrukce má přímou vazbu na sociální dialog: nárůst podílu \ac{OSVČ} na zaměstnanosti (srov.~obr.~\ref{fig:self_employment_map}), částečně podnícený vyhýbáním se kolektivnímu vyjednávání ze strany zaměstnavatelů (tzv.\ švarcsystém), oslabuje příjmovou základnu důchodového systému.
Slabé kolektivní vyjednávání tak působí systémově na dvou úrovních: přímou cestou stlačuje mzdy zaměstnanců a~nepřímo, prostřednictvím rozšiřování fiktivní samostatné výdělečné činnosti, erozi pojistné základny, z~níž se důchody financují.
Reforma sociálního dialogu --- zejména posílení mechanismů proti švarcsystému a~rozšíření sektorových \ac{KS} --- by přispěla i~ke stabilitě průběžného důchodového systému.

Parametrický výpočet starobního důchodu v~závislosti na úrovni příjmů odhaluje zásadní asymetrii mezi skupinami pojištěnců: zaměstnanec s~průměrnou mzdou dosáhne výrazně vyššího důchodu než \ac{OSVČ} se srovnatelnými celkovými náklady zaměstnavatele, a~to zejména v~případě paušálního daňového režimu, kde je výměra důchodu stropována pevnou platbou.
Tato konstrukce má přímou vazbu na sociální dialog: nárůst podílu \ac{OSVČ} na zaměstnanosti (srov.~obr.~\ref{fig:self_employment_map}), částečně podnícený vyhýbáním se kolektivnímu vyjednávání ze strany zaměstnavatelů (tzv.\ švarcsystém), oslabuje příjmovou základnu důchodového systému.
Slabé kolektivní vyjednávání tak působí systémově na dvou úrovních: přímou cestou stlačuje mzdy zaměstnanců a~nepřímo, prostřednictvím rozšiřování fiktivní samostatné výdělečné činnosti, erozi pojistné základny, z~níž se důchody financují.
Reforma sociálního dialogu --- zejména posílení mechanismů proti švarcsystému a~rozšíření sektorových \ac{KS} --- by přispěla i~ke stabilitě průběžného důchodového systému.

Parametrický výpočet starobního důchodu v~závislosti na úrovni příjmů odhaluje zásadní asymetrii mezi skupinami pojištěnců: zaměstnanec s~průměrnou mzdou dosáhne výrazně vyššího důchodu než \ac{OSVČ} se srovnatelnými celkovými náklady zaměstnavatele, a~to zejména v~případě paušálního daňového režimu, kde je výměra důchodu stropována pevnou platbou.
Tato konstrukce má přímou vazbu na sociální dialog: nárůst podílu \ac{OSVČ} na zaměstnanosti (srov.~obr.~\ref{fig:self_employment_map}), částečně podnícený vyhýbáním se kolektivnímu vyjednávání ze strany zaměstnavatelů (tzv.\ švarcsystém), oslabuje příjmovou základnu důchodového systému.
Slabé kolektivní vyjednávání tak působí systémově na dvou úrovních: přímou cestou stlačuje mzdy zaměstnanců a~nepřímo, prostřednictvím rozšiřování fiktivní samostatné výdělečné činnosti, erozi pojistné základny, z~níž se důchody financují.
Reforma sociálního dialogu --- zejména posílení mechanismů proti švarcsystému a~rozšíření sektorových \ac{KS} --- by přispěla i~ke stabilitě průběžného důchodového systému.

\input{texparts/python/cz_pension_income} Osa x odpovídá celkovým výdajům plátce: pro zaměstnance zahrnuje hrubou mzdu i odvody zaměstnavatele (\SI{33.8}{\percent}); pro \ac{OSVČ} se standardními odvody je to zisk (příjmy\,$-$\,výdaje); pro \ac{OSVČ} v~paušálním daňovém režimu jde o~měsíční příjmy (revenue), přičemž výše důchodu je v~každém pásmu pevná. Parametry roku~2026, předpokládaná pojistná doba 40~let. Výpočet dle zákona č.\,155/1995~Sb.~\cite{zakon_zpds_1995} (zákon o~důchodovém pojištění), zákona č.\,270/2023~Sb.~\cite{zakon_duchreforma_2023} (důchodová reforma) a nařízení vlády č.\,365/2025~Sb.~\cite{nv_365_2025}.
 Osa x odpovídá celkovým výdajům plátce: pro zaměstnance zahrnuje hrubou mzdu i odvody zaměstnavatele (\SI{33.8}{\percent}); pro \ac{OSVČ} se standardními odvody je to zisk (příjmy\,$-$\,výdaje); pro \ac{OSVČ} v~paušálním daňovém režimu jde o~měsíční příjmy (revenue), přičemž výše důchodu je v~každém pásmu pevná. Parametry roku~2026, předpokládaná pojistná doba 40~let. Výpočet dle zákona č.\,155/1995~Sb.~\cite{zakon_zpds_1995} (zákon o~důchodovém pojištění), zákona č.\,270/2023~Sb.~\cite{zakon_duchreforma_2023} (důchodová reforma) a nařízení vlády č.\,365/2025~Sb.~\cite{nv_365_2025}.
 Osa x odpovídá celkovým výdajům plátce: pro zaměstnance zahrnuje hrubou mzdu i odvody zaměstnavatele (\SI{33.8}{\percent}); pro \ac{OSVČ} se standardními odvody je to zisk (příjmy\,$-$\,výdaje); pro \ac{OSVČ} v~paušálním daňovém režimu jde o~měsíční příjmy (revenue), přičemž výše důchodu je v~každém pásmu pevná. Parametry roku~2026, předpokládaná pojistná doba 40~let. Výpočet dle zákona č.\,155/1995~Sb.~\cite{zakon_zpds_1995} (zákon o~důchodovém pojištění), zákona č.\,270/2023~Sb.~\cite{zakon_duchreforma_2023} (důchodová reforma) a nařízení vlády č.\,365/2025~Sb.~\cite{nv_365_2025}.
 Osa x odpovídá celkovým výdajům plátce: pro zaměstnance zahrnuje hrubou mzdu i odvody zaměstnavatele (\SI{33.8}{\percent}); pro \ac{OSVČ} se standardními odvody je to zisk (příjmy\,$-$\,výdaje); pro \ac{OSVČ} v~paušálním daňovém režimu jde o~měsíční příjmy (revenue), přičemž výše důchodu je v~každém pásmu pevná. Parametry roku~2026, předpokládaná pojistná doba 40~let. Výpočet dle zákona č.\,155/1995~Sb.~\cite{zakon_zpds_1995} (zákon o~důchodovém pojištění), zákona č.\,270/2023~Sb.~\cite{zakon_duchreforma_2023} (důchodová reforma) a nařízení vlády č.\,365/2025~Sb.~\cite{nv_365_2025}.
